\subsection{Bisection Method}

\begin{frame}{Bisection Method: Explicación Gráfica}
    \begin{block}{Concepto principal}
        [Aguilar: Explica cómo en optimización usamos bisección para buscar la raíz de la primera derivada $f'(x)=0$].
    \end{block}
    \begin{itemize}
        \item {[Idea gráfica 1]}
        \item {[Idea gráfica 2]}
    \end{itemize}
\end{frame}

\begin{frame}{Bisection Method: Pseudocódigo}
    \begin{algorithmic}[1]
        \State \textbf{Input:} [La derivada $f'$, límites $a$ y $b$]
        \State [Lógica evaluando signos de la derivada]
        \State \textbf{Return} [Resultado]
    \end{algorithmic}
\end{frame}

\begin{frame}{Bisection Method: Ejemplo Paso a Paso}
    \begin{exampleblock}{Planteamiento del problema}
        [Aguilar: Función de prueba y cómo evaluamos su derivada].
    \end{exampleblock}
    \begin{itemize}
        \item \textbf{Iteración 1:} [Evaluación en el punto medio]
        \item \textbf{Iteración 2:} [Corte del intervalo]
    \end{itemize}
\end{frame}

\begin{frame}[fragile]{Bisection Method: Código}
\begin{lstlisting}[language=Python]
# [Aguilar: Reemplaza esto con tu código real]
def bisection_ejemplo():
    pass
\end{lstlisting}
\end{frame}

\begin{frame}{Bisection Method: Análisis de Convergencia}
    \begin{alertblock}{Tasa de Convergencia}
        [Aguilar: Habla sobre la reducción exacta a la mitad y convergencia lineal].
    \end{alertblock}
    \begin{itemize}
        \item {[Requisito estricto de derivabilidad de la función]}.
    \end{itemize}
\end{frame}